\documentclass[11pt,oneside]{article}
\usepackage[a4paper,margin=27mm,headheight=15pt]{geometry}
\usepackage[T1]{fontenc}
\usepackage[utf8]{inputenc}
\IfFileExists{libertinus.sty}{\usepackage{libertinus}}{\usepackage{lmodern}}
\usepackage{microtype}
\usepackage{amsmath,amssymb,amsthm,mathtools,mathrsfs,bm}
\usepackage{booktabs,longtable,array,tabularx,multirow}
\usepackage{xcolor}
\usepackage{enumitem}
\usepackage{fancyhdr}
\usepackage{titlesec}
\usepackage{xurl}
\usepackage{hyperref}
\usepackage[nameinlink,noabbrev]{cleveref}
\definecolor{linkblue}{RGB}{25,75,135}
\definecolor{shadegray}{RGB}{246,247,249}
\definecolor{rulegray}{RGB}{195,201,208}
\hypersetup{
  colorlinks=true,
  linkcolor=linkblue,
  citecolor=linkblue,
  urlcolor=linkblue,
  pdftitle={Binary Spectral-Endpoint Bridges for the Fabius and Rvachev Functions},
  pdfsubject={q-Pochhammer zero divisors, Thue-Morse endpoint weights, Mellin strip duality, and branch-resolved Lambert maps},
  pdfkeywords={Fabius function, Rvachev up-function, Thue-Morse, q-binomial, q-Pochhammer, sinc product, Lambert W, Mellin transform}
}
\pagestyle{fancy}
\fancyhf{}
\fancyhead[L]{\small Fabius--Rvachev frontier report}
\fancyhead[R]{\small binary $q$--TM--sinc--Lambert bridges}
\fancyfoot[C]{\thepage}
\renewcommand{\headrulewidth}{0.4pt}
\titleformat{\section}{\Large\bfseries}{\thesection}{0.7em}{}
\titleformat{\subsection}{\large\bfseries}{\thesubsection}{0.7em}{}
\titleformat{\subsubsection}{\normalsize\bfseries}{\thesubsubsection}{0.7em}{}
\setcounter{secnumdepth}{3}
\setcounter{tocdepth}{2}
\setlist[itemize]{leftmargin=2em,itemsep=0.25em,topsep=0.4em}
\setlist[enumerate]{leftmargin=2.3em,itemsep=0.25em,topsep=0.4em}
\setlength{\emergencystretch}{2em}
\newtheorem{theorem}{Theorem}[section]
\newtheorem{proposition}[theorem]{Proposition}
\newtheorem{lemma}[theorem]{Lemma}
\newtheorem{corollary}[theorem]{Corollary}
\newtheorem{conjecture}[theorem]{Conjecture}
\theoremstyle{definition}
\newtheorem{definition}[theorem]{Definition}
\newtheorem{algorithm}[theorem]{Algorithm}
\newtheorem{example}[theorem]{Example}
\theoremstyle{remark}
\newtheorem{remark}[theorem]{Remark}
\newtheorem{warning}[theorem]{Warning}
\newcommand{\R}{\mathbb R}
\newcommand{\C}{\mathbb C}
\newcommand{\Q}{\mathbb Q}
\newcommand{\N}{\mathbb N}
\newcommand{\Z}{\mathbb Z}
\newcommand{\E}{\mathbb E}
\newcommand{\Prob}{\mathbb P}
\newcommand{\Var}{\operatorname{Var}}
\newcommand{\Up}{\operatorname{up}}
\newcommand{\sinc}{\operatorname{sinc}}
\newcommand{\e}{\mathrm e}
\newcommand{\ii}{\mathrm i}
\newcommand{\dd}{\,\mathrm d}
\newcommand{\bigO}{\mathcal O}
\newcommand{\smallo}{o}
\newcommand{\supp}{\operatorname{supp}}
\newcommand{\ord}{\operatorname{ord}}
\newcommand{\Res}{\operatorname*{Res}}
\newcommand{\vtwo}{v_2}
\newcommand{\wt}{\operatorname{wt}_2}
\newcommand{\Mellin}{\mathcal M}
\newcommand{\falling}[2]{(#1)_{\underline{#2}}}
\newcommand{\qbinom}[3]{\genfrac{[}{]}{0pt}{}{#1}{#2}_{#3}}
\newcommand{\statusE}{\textsf{[E]}}
\newcommand{\statusN}{\textsf{[N]}}
\newcommand{\statusC}{\textsf{[C]}}
\newenvironment{boxedremark}{%
  \par\smallskip\noindent\begin{center}\begin{minipage}{0.94\textwidth}%
  \color{black}\hrule\vspace{0.6em}\small
}{%
  \vspace{0.6em}\hrule\end{minipage}\end{center}\smallskip
}
\newenvironment{proofidea}{\par\smallskip\noindent\textit{Proof architecture.}\ }{\hfill$\square$\par\smallskip}

\title{\bfseries Binary Spectral--Endpoint Bridges for the Fabius and Rvachev Functions\\[0.35em]
\large $q$-Pochhammer zero divisors, Thue--Morse endpoint weights, Mellin strip duality, and branch-resolved Lambert maps}
\author{Frontier research report based on a recursive audit of\\
\small Vladimir Reshetnikov's \texttt{ProveIt/Analysis/FabiusFunction/docs} corpus}
\date{27 August 2026}

\begin{document}
\maketitle

\begin{abstract}
This report audits the recursive LaTeX corpus under
\path{Analysis/FabiusFunction/docs} and then develops a new synthesis joining five strands
that the existing documents largely treat in parallel: the Thue--Morse Mahler product,
base-$1/2$ and base-$1/4$ $q$-calculus, the Rvachev infinite sinc product, the exact
negative-Laplace decomposition at the Fabius endpoint, and the lower real branch of
Lambert's $W$ function.  The audit covers 57 \texttt{.tex} files: five living synthesis
or reference documents, thirteen active frontier drafts, thirty-two frozen historical
versions, and seven transcribed or corrected source papers.  Frozen duplicate snapshots
were checked through the archive's provenance map and through the living documents into
which their distinctive material was merged.

The principal new results, in the precise sense ``not found explicitly in the inspected
corpus'', are as follows.  First, the Rvachev transform admits the $q$-Pochhammer product
\[
 \widehat\Up(z)=\prod_{\ell\ge1}(z^2/\ell^2;1/4)_\infty
               =\prod_{n\ge1}(1-z^2/n^2)^{\vtwo(n)+1},
\]
so its integer zero multiplicity is exactly $\vtwo(n)+1$.  A finite version has both a
Gaussian-binomial factorization and an exact signed Thue--Morse exponential-sum
factorization.  Second, if
$\Theta(q)=\sum_{n\ge0}(-1)^{\wt(n)}q^n=\prod_{j\ge0}(1-q^{2^j})$, then the exponentially
small endpoint correction in the negative-Laplace formula is exactly
$-\log\Theta(\e^{-s})$.  Its Lambert-series coefficients satisfy
\[
 -\log\Theta(\e^{-s})
 =\sum_{n\ge1}\frac{2^{\vtwo(n)+1}-1}{n}\e^{-ns}
 =\sum_{n\ge1}\frac{2^{\ord_{z=n}\widehat\Up}-1}{n}\e^{-ns},
\]
giving a direct spectral--endpoint dictionary.  Third, its Mellin transform is
\[
 \int_0^\infty[-\log\Theta(\e^{-s})]s^{w-1}\dd s
 =\frac{\Gamma(w)\zeta(w+1)}{1-2^{-w}},
\]
which realizes, on the complementary positive Mellin strip, the same meromorphic
$\Gamma\zeta$ kernel used by the corpus for the negative-Laplace logarithm.  This gives a
new Thue--Morse/weighted-zero derivation of every Fourier coefficient of the endpoint's
log-periodic function and recovers the corpus's already established exact mean in terms of
the first Stieltjes constant.  Fourth, the exact endpoint saddle is the fixed point of a
periodically perturbed Lambert map, with the lower branch governing the small-argument
endpoint regime and an explicit derivative proving local contraction there.  Finally, an
all-orders base-$1/4$ tail expansion places the corpus's established centered moments in
spectral tail coordinates and yields a certified acceleration of finite Thue--Morse
approximants to the Fourier transform.

No assertion of global literature priority is made.  Every result is labeled as established
in the source corpus, newly proved here relative to that corpus, or conjectural/future work.
\end{abstract}

\tableofcontents
\newpage

\section{Scope, audit method, and proof-status discipline}

\subsection{What was inspected}
The repository separates formalization-backed exposition from non-formalized frontier
mathematics.  The living primary article records only statements with named Lean
counterparts, whereas the frontier volume deliberately quarantines plausible or proved-on-
paper results not yet represented by an exact formal declaration.  The recursive source
audit used that distinction rather than flattening all files into one undifferentiated
literature review.

At the snapshot date, the recursive corpus contains the following groups.
\begin{center}
\small
\begin{tabularx}{\textwidth}{@{}l c >{\raggedright\arraybackslash}X@{}}
\toprule
Group & Files & Role and audit treatment\\
\midrule
Living documents & 5 & Primary exposition, glossary, non-elementarity, Lean walkthrough,
and frontier synthesis; used as the theorem and notation baseline.\\
Active drafts & 13 & Thue--Morse atlas and twelve Fourier-decay studies; checked for
newer claims, conflicts, and unresolved questions.\\
Frozen archive & 32 & First versions and historical parallel drafts; checked through the
archive provenance descriptions and merged-content map, with distinctive claims traced into
living successors.\\
Source papers & 7 & Corrected or transcribed mathematical papers; used for provenance,
comparison, and terminology.\\
\midrule
Total & 57 & Recursive LaTeX corpus audited.\\
\bottomrule
\end{tabularx}
\end{center}
The full path-level inventory is recorded in \cref{app:inventory}.

\subsection{Status labels}
Throughout the report:
\begin{itemize}
\item \statusE\ means an established theorem or audited conclusion already present in the
repository corpus (often formally verified in Lean);
\item \statusN\ means a theorem proved in this report and not located explicitly in the
inspected corpus by path, phrase, and formula searches;
\item \statusC\ means a conjecture, program, or interpretation whose proof is not supplied.
\end{itemize}
``New'' is therefore a corpus-relative statement, not an unsupported claim of first
publication in the worldwide literature.

A final formula-by-formula overlap audit corrected three preliminary classifications before
release.  The exact mean and full Fourier spectrum of the endpoint-periodic function are
already established in the corpus; so are the centered moment generating function, its
initial rational moments, and a Bernoulli recurrence for those moments.  Moreover, the
same meromorphic $\Gamma\zeta$ kernel already occurs there as the Mellin transform of
$\Lambda$ on $-1<\Re w<0$.  The new claims below are therefore restricted to the
$q$-Pochhammer and weighted-zero dictionaries, the complementary positive-strip
realization of that kernel, the resulting Thue--Morse/zero-divisor derivations, the exact
Lambert branch map, and the certified finite-product consequences.

\subsection{The structural gap that motivates the report}
The corpus already develops, separately and deeply,
\begin{enumerate}
\item exact dyadic values of the Fabius function involving $(1/2;1/2)_n$ and Gaussian
binomial coefficients;
\item the Thue--Morse product $\prod_{j\ge0}(1-q^{2^j})$ and its Mahler equation;
\item the Rvachev Fourier transform as an infinite product of geometrically rescaled sinc
functions;
\item a quadratic-plus-periodic negative-Laplace decomposition and a lower-Lambert endpoint
phase;
\item a multi-regime theory of Fourier decay.
\end{enumerate}
What is missing is an exact dictionary showing that the same binary arithmetic controls
\emph{both} the zero divisor of the Fourier transform and the exponentially small correction
in the endpoint Laplace transform.  Establishing that dictionary leads to a new
$q$-Pochhammer zero-divisor formula, a weighted-zero interpretation of the corpus's
$\Gamma\zeta$ Mellin kernel, exact tilted-cumulant recurrences, and branch-resolved
Lambert-$W$ maps.

\begin{boxedremark}
\textbf{Main synthesis in one line.}
The Thue--Morse product controls the endpoint correction; its logarithmic coefficients are
$2^{\mu(n)}-1$, where $\mu(n)$ is exactly the multiplicity of the Fourier zero at the integer
$n$.  Mellin transformation turns the resulting $2$-adic Dirichlet series into the
$\Gamma\zeta$ Fourier modes of the endpoint periodic function, and that periodic function
perturbs the lower-Lambert saddle.
\end{boxedremark}

\section{Established framework and normalization}

\subsection{The three functions}
\statusE\ The corpus carefully distinguishes three related global objects.
The bounded Fabius function $F$ is the distribution function on $[0,1]$ satisfying
\[
 F'(x)=2F(2x)\quad(0\le x\le\tfrac12),\qquad
 F(1-x)=1-F(x),
\]
with constant extensions outside $[0,1]$.  Rvachev's even compactly supported bump is
\[
 \Up(x)=F(1-|x|)\quad(|x|\le1),\qquad \Up(x)=0\quad(|x|\ge1).
\]
The signed global extension $\mathcal F$ satisfies the pure dilation equation on all of
$\R$ and should not be silently substituted for $F$ or $\Up$.

\subsection{Probability model}
Let $U_1,U_2,\ldots$ be independent $\operatorname{Unif}[0,1]$ variables and set
\begin{equation}
 X=\sum_{j\ge1}2^{-j}U_j.                                  \label{eq:X}
\end{equation}
\statusE\ Then $F(x)=\Prob(X\le x)$.  The centered variable
\begin{equation}
 Y=2X-1=\sum_{n\ge0}2^{-n}V_n,
 \qquad V_n\overset{\mathrm{iid}}\sim\operatorname{Unif}[-1/2,1/2], \label{eq:Y}
\end{equation}
has density $\Up$.  Consequently $Y$ is symmetric, $|Y|\le1$, and its Fourier transform
under the convention
\[
 \widehat f(z)=\int_\R f(x)\e^{-2\pi\ii xz}\dd x
\]
is a characteristic function on the real axis.

\subsection{Sinc product}
Write $\sinc w=\sin w/w$ with the removable value $1$ at $w=0$, and set
\begin{equation}
 \Phi(z)=\widehat\Up(z).
\end{equation}
\statusE\ The exact product and refinement law are
\begin{equation}
 \boxed{\Phi(z)=\prod_{n=0}^{\infty}\sinc\!\left(\frac{\pi z}{2^n}\right),
 \qquad \Phi(z)=\sinc(\pi z)\Phi(z/2).}                    \label{eq:Phi-sinc}
\end{equation}
The product is entire and also admits the weighted cosine factorization
$\Phi(z)=\prod_{m\ge1}\cos(\pi z/2^m)^m$.

\subsection{Thue--Morse product}
Let
\begin{equation}
 \varepsilon_n=(-1)^{\wt(n)},\qquad
 \Theta(q)=\sum_{n\ge0}\varepsilon_n q^n.                 \label{eq:Theta-def}
\end{equation}
\statusE\ For $|q|<1$,
\begin{equation}
 \boxed{\Theta(q)=\prod_{j\ge0}(1-q^{2^j}),
 \qquad \Theta(q)=(1-q)\Theta(q^2).}                       \label{eq:Theta-product}
\end{equation}
At finite depth,
\begin{equation}
 \sum_{k=0}^{2^m-1}\varepsilon_k q^k
 =\prod_{j=0}^{m-1}(1-q^{2^j}),                            \label{eq:TM-finite}
\end{equation}
and Prouhet cancellation gives
\begin{equation}
 \sum_{k=0}^{2^m-1}\varepsilon_k k^r=0\ (0\le r<m),
 \qquad
 \sum_{k=0}^{2^m-1}\varepsilon_k k^m=(-1)^m2^{\binom m2}m!. \label{eq:Prouhet}
\end{equation}

\subsection{Negative Laplace transform and periodic correction}
Define
\begin{equation}
 B(s)=\E\e^{-sX}
 =\prod_{j\ge1}\frac{1-\e^{-s/2^j}}{s/2^j},
 \qquad \Lambda(s)=\log B(s),\quad s>0.                   \label{eq:B}
\end{equation}
\statusE\ The exact dilation relation is
\begin{equation}
 \Lambda(2s)=\Lambda(s)+\log(1-\e^{-s})-\log s.           \label{eq:Lambda-dilation}
\end{equation}
Let $L=\log2$ and
\begin{equation}
 T(s)=\sum_{m\ge0}\log(1-\e^{-2^ms}),\qquad
 \mathcal E(s)=-T(s)>0.                                   \label{eq:E-tail}
\end{equation}
Then
\begin{equation}
 P(u)=\Lambda(2^u)+\frac L2(u^2-u)+T(2^u)                 \label{eq:P-def}
\end{equation}
is smooth and $1$-periodic, and
\begin{equation}
 \boxed{\Lambda(s)=-\frac{(\log s)^2}{2L}+\frac12\log s
 +P(\log_2s)+\mathcal E(s).}                              \label{eq:Lambda-periodic}
\end{equation}
The corpus gives every Fourier mode of $P$, including its exact mean, and proves that no
nonzero mode vanishes.  It also obtains
$\Gamma(w)\zeta(w+1)/(1-2^{-w})$ as the Mellin transform of $\Lambda$ on
$-1<\Re w<0$.  One of our goals is to realize that same meromorphic kernel from the
Thue--Morse endpoint tail on the complementary strip $\Re w>0$, and thereby rederive the
entire periodic spectrum by one weighted-zero residue calculation.

\subsection{The Fourier-decay spectrum is not one number}
\statusE\ The comparative Fourier-decay audits establish a common lognormal core
\[
 \exp\!\left[-\frac{(\log x)^2}{2\log2}\right]
\]
but different power corrections for different notions of size.  In the notation of the
audits, the ordering is
\[
 \kappa_\infty<\kappa_2<\kappa_1<\kappa_0,
\]
corresponding respectively to extremal rays, quadratic mean, Ces\`aro mean, and typical or
almost-everywhere behavior.  The second-wave audit additionally verifies that an
unrestricted Ces\`aro normalization can be attained by a positive real-analytic decreasing
gauge, although every stabilizing or bounded-distortion gauge fails.  This report does not
replace that multi-regime theory; instead, it supplies exact arithmetic and transform
identities that sit underneath all regimes.

\section{A geometric $q$-Rvachev family}
\label{sec:qfamily}

\subsection{Definition and probabilistic realization}
The binary Rvachev product belongs naturally to a one-parameter family.

\begin{definition}[$q$-Rvachev transform]
For $0<q<1$, define
\begin{equation}
 \Phi_q(z)=\prod_{n\ge0}\sinc(\pi q^{n/2}z).               \label{eq:Phiq}
\end{equation}
Let $V_0,V_1,\ldots$ be independent $\operatorname{Unif}[-1/2,1/2]$ variables and put
\begin{equation}
 Y_q=\sum_{n\ge0}q^{n/2}V_n.                               \label{eq:Yq}
\end{equation}
\end{definition}

\begin{theorem}[Probability, support, smooth density, and refinement]\label{thm:q-density}
\statusN\ The product \eqref{eq:Phiq} converges locally uniformly to an entire function.
It is the Fourier transform of the law of $Y_q$, whose support is
\[
 \left[-\frac1{2(1-\sqrt q)},\frac1{2(1-\sqrt q)}\right].
\]
The law has a compactly supported $C^\infty$ density $u_q$, and
\begin{equation}
 \Phi_q(z)=\sinc(\pi z)\Phi_q(\sqrt q\,z).                 \label{eq:q-refine-Fourier}
\end{equation}
Writing $a=\sqrt q$, the density satisfies
\begin{align}
 u_q(x)&=\int_{(x-1/2)/a}^{(x+1/2)/a}u_q(y)\dd y,           \label{eq:q-density-integral}\\
 u_q'(x)&=\frac1a\left[u_q\!\left(\frac{x+1/2}{a}\right)
                     -u_q\!\left(\frac{x-1/2}{a}\right)\right]. \label{eq:q-density-diff}
\end{align}
At $q=1/4$, $Y_q=Y$, $u_q=\Up$, and \eqref{eq:q-density-diff} is the usual Rvachev
functional-differential equation.
\end{theorem}

\begin{proof}
The summands in \eqref{eq:Yq} converge absolutely almost surely and are bounded by a
geometric series.  Independence gives \eqref{eq:Phiq}.  Since
$\sinc(\pi q^{n/2}z)=1+\bigO(q^n|z|^2)$ locally uniformly, the entire product converges on
compact sets.  If $N\sim\log|x|/|\log\sqrt q|$, bounding the first $N$ sinc factors by a
constant times $(|x|q^{n/2})^{-1}$ gives
\[
 \log|\Phi_q(x)|\le-\frac{(\log|x|)^2}{2|\log\sqrt q|}+\bigO(\log|x|),
\]
so every polynomial multiple is integrable.  Fourier inversion therefore yields a
$C^\infty$ density.  Finally $Y_q\overset d=V_0+aY_q'$ with an independent copy $Y_q'$.
Convolving with the uniform density gives \eqref{eq:q-density-integral}; differentiation
gives \eqref{eq:q-density-diff}.
\end{proof}

\subsection{$q$-Pochhammer factorization}
Recall
\[
 (a;q)_m=\prod_{j=0}^{m-1}(1-aq^j),\qquad
 (a;q)_\infty=\prod_{j\ge0}(1-aq^j).
\]

\begin{theorem}[Euler--$q$ factorization]\label{thm:qpoch-sinc}
\statusN\ For every $0<q<1$ and $z\in\C$,
\begin{align}
 \Phi_q(z)&=\prod_{\ell=1}^{\infty}(z^2/\ell^2;q)_\infty,  \label{eq:qpoch-full}\\
 \Phi_{q,m}(z):=\prod_{n=0}^{m-1}\sinc(\pi q^{n/2}z)
 &=\prod_{\ell=1}^{\infty}(z^2/\ell^2;q)_m.                \label{eq:qpoch-finite}
\end{align}
For $|z|<1$,
\begin{align}
 \log\Phi_q(z)
 &=-\sum_{r\ge1}\frac{\zeta(2r)}{r(1-q^r)}z^{2r},         \label{eq:q-log-full}\\
 \log\Phi_{q,m}(z)
 &=-\sum_{r\ge1}\frac{1-q^{mr}}{1-q^r}\frac{\zeta(2r)}r z^{2r}. \label{eq:q-log-finite}
\end{align}
\end{theorem}

\begin{proof}
Euler's sine product gives
\[
 \sinc(\pi q^{n/2}z)=\prod_{\ell\ge1}\left(1-\frac{q^nz^2}{\ell^2}\right).
\]
The double product is absolutely and locally uniformly convergent because
$\sum_{n,\ell}q^n/\ell^2<\infty$, so the factors may be regrouped by $\ell$.  This proves
\eqref{eq:qpoch-full} and \eqref{eq:qpoch-finite}.  Expanding
$\log(1-w)=-\sum_{r\ge1}w^r/r$ and summing the geometric progression in $n$ proves the
logarithmic formulas.
\end{proof}

\begin{corollary}[Gaussian-binomial finite factorization]\label{cor:qbinom-finite}
\statusN\ The finite product has the exact representation
\begin{equation}
 \boxed{
 \Phi_{q,m}(z)=\prod_{\ell\ge1}
 \left[
  \sum_{j=0}^{m}(-1)^jq^{\binom j2}\qbinom mjq\frac{z^{2j}}{\ell^{2j}}
 \right].}                                                  \label{eq:qbinom-sinc}
\end{equation}
\end{corollary}

\begin{proof}
Apply the finite $q$-binomial theorem
$(a;q)_m=\sum_{j=0}^m(-1)^jq^{\binom j2}\qbinom mjq a^j$
to each factor in \eqref{eq:qpoch-finite}.
\end{proof}

\begin{remark}
The living corpus uses $q=1/2$ Gaussian coefficients in exact dyadic values of $F$.
Equation \eqref{eq:qbinom-sinc} reveals a second, independent $q$-geometry: the spectral
product naturally uses $q=1/4$, the square of the spatial dilation.  The two bases should
not be conflated, but their coexistence is structural rather than accidental.
\end{remark}

\subsection{Cumulants in the $q$-family}
Let $c_N(q)=\E Y_q^{2N}$, with $c_0(q)=1$, and use the standard Bernoulli numbers
$B_{2r}$.

\begin{theorem}[$q$-Bernoulli cumulants and moments]\label{thm:q-cumulants}
\statusN\ All odd cumulants of $Y_q$ vanish, and
\begin{equation}
 \boxed{\kappa_{2r}(Y_q)=\frac{B_{2r}}{2r(1-q^r)}.}         \label{eq:q-cumulant}
\end{equation}
The even moments obey
\begin{equation}
 \boxed{
 c_N(q)=\frac1{2N}\sum_{r=1}^{N}
 \binom{2N}{2r}\frac{B_{2r}}{1-q^r}\,c_{N-r}(q).}         \label{eq:q-moment-rec}
\end{equation}
\end{theorem}

\begin{proof}
Insert $z=t/(2\pi)$ in \eqref{eq:q-log-full} and compare
$\log\E\e^{\ii tY_q}$ with
$\sum_{n\ge1}\kappa_n(\ii t)^n/n!$.  Euler's evaluation of $\zeta(2r)$ gives
\eqref{eq:q-cumulant}.  The standard moment--cumulant recursion
$m_n=\sum_{k=1}^n\binom{n-1}{k-1}\kappa_km_{n-k}$, followed by
$\binom{2N-1}{2r-1}=\frac rN\binom{2N}{2r}$, gives
\eqref{eq:q-moment-rec}.
\end{proof}

\section{The binary specialization: a $2$-adic zero divisor}
\label{sec:zero-divisor}

Set $q=1/4$.  Then $\Phi_{1/4}=\Phi=\widehat\Up$.

\begin{theorem}[Arithmetic canonical product]\label{thm:arithmetic-product}
\statusN\ The Rvachev transform has the genus-zero product
\begin{equation}
 \boxed{
 \Phi(z)=\prod_{n=1}^{\infty}\left(1-\frac{z^2}{n^2}\right)^{\mu(n)},
 \qquad \mu(n)=\vtwo(n)+1.}                                \label{eq:arithmetic-product}
\end{equation}
Thus the nonzero zeros of $\Phi$ are exactly the integers, and
\begin{equation}
 \boxed{\ord_{z=n}\Phi=\vtwo(|n|)+1\qquad(n\in\Z\setminus\{0\}).} \label{eq:zero-mult}
\end{equation}
Moreover
\begin{equation}
 \sum_{n=1}^{N}\mu(n)=2N-\wt(N).                           \label{eq:zero-count}
\end{equation}
\end{theorem}

\begin{proof}
From \eqref{eq:qpoch-full},
\[
 \Phi(z)=\prod_{j\ge0}\prod_{k\ge1}
 \left(1-\frac{z^2}{4^jk^2}\right).
\]
For a fixed positive integer $n$, the equality $n=2^jk$ with $k\in\N$ occurs exactly for
$j=0,\ldots,\vtwo(n)$.  Absolute convergence permits regrouping and proves
\eqref{eq:arithmetic-product}.  Each contributing factor has a simple zero, proving
\eqref{eq:zero-mult}.  Finally
\[
 \sum_{n\le N}\mu(n)=N+\sum_{n\le N}\vtwo(n)
 =N+\vtwo(N!)=2N-\wt(N)
\]
by Legendre's formula.
\end{proof}

\begin{corollary}[Dirichlet series of the zero multiplicities]\label{cor:mu-Dirichlet}
For $\Re w>1$,
\begin{equation}
 \boxed{\sum_{n\ge1}\frac{\mu(n)}{n^w}
 =\frac{\zeta(w)}{1-2^{-w}}.}                              \label{eq:mu-Dirichlet}
\end{equation}
Consequently \eqref{eq:q-log-full} at $q=1/4$ can be read directly from the zero divisor:
\begin{equation}
 \log\Phi(z)=-\sum_{r\ge1}\frac{z^{2r}}r
 \sum_{n\ge1}\frac{\mu(n)}{n^{2r}}
 =-\sum_{r\ge1}\frac{4^r\zeta(2r)}{r(4^r-1)}z^{2r}.      \label{eq:logPhi-mu}
\end{equation}
\end{corollary}

\begin{proof}
Since $\mu(n)=\sum_{j\ge0}\mathbf1_{2^j\mid n}$, Tonelli's theorem gives
\[
 \sum_n\mu(n)n^{-w}=\sum_{j\ge0}\sum_{2^j\mid n}n^{-w}
 =\zeta(w)\sum_{j\ge0}2^{-jw}.
\]
\end{proof}

\subsection{Finite triple factorization: sinc, $q$-Pochhammer, and Thue--Morse}
Define
\begin{equation}
 \Phi_m(z)=\prod_{n=0}^{m-1}\sinc\!\left(\frac{\pi z}{2^n}\right),
 \qquad
 S_m(z)=\sum_{k=0}^{2^m-1}\varepsilon_k
 \e^{-2\pi\ii kz/2^{m-1}}.                                \label{eq:Phi-m-Sm}
\end{equation}

\begin{theorem}[Finite triple factorization]\label{thm:triple-factor}
\statusN\ For $m\ge1$,
\begin{equation}
 \boxed{
 \prod_{\ell\ge1}(z^2/\ell^2;1/4)_m
 =\Phi_m(z)
 =\frac{2^{m(m-3)/2}\e^{2\pi\ii z(1-2^{-m})}}{(\ii\pi z)^m}
 S_m(z).}                                                   \label{eq:triple-factor}
\end{equation}
The apparent singularity at $z=0$ is removable.  In fact $S_m$ has a zero of exact order
$m$ there, with
\begin{equation}
 S_m(z)=(2\pi\ii z)^m2^{-m(m-1)/2}+\bigO(z^{m+1}).         \label{eq:Sm-leading}
\end{equation}
\end{theorem}

\begin{proof}
The first equality is \eqref{eq:qpoch-finite}.  Use
$\sin w=\e^{\ii w}(1-\e^{-2\ii w})/(2\ii)$ in each sinc factor.  The powers of two give
$2^{\sum_{n<m}n-m}=2^{m(m-3)/2}$ and the phases sum to
$\e^{2\pi\ii z(1-2^{-m})}$.  The remaining product is
\[
 \prod_{n=0}^{m-1}(1-\e^{-2\pi\ii z/2^n})
 =\prod_{j=0}^{m-1}(1-w^{2^j})=\sum_{k<2^m}\varepsilon_kw^k,
 \quad w=\e^{-2\pi\ii z/2^{m-1}},
\]
which is $S_m(z)$.  Equation \eqref{eq:Prouhet} annihilates the first $m$ Taylor
coefficients of $S_m$ and gives the leading coefficient in \eqref{eq:Sm-leading}.
\end{proof}

\begin{corollary}[A Gaussian-binomial identity for a Thue--Morse exponential sum]
Combining \eqref{eq:qbinom-sinc} and \eqref{eq:triple-factor} gives
\begin{align}
 S_m(z)
 &=(\ii\pi z)^m2^{-m(m-3)/2}\e^{-2\pi\ii z(1-2^{-m})} \notag\\
 &\quad\times
 \prod_{\ell\ge1}
 \left[
  \sum_{j=0}^{m}(-1)^j4^{-\binom j2}\qbinom mj{1/4}
  \frac{z^{2j}}{\ell^{2j}}
 \right].                                                  \label{eq:TM-qbinom}
\end{align}
This is an exact bridge between a finite signed Thue--Morse block and Gaussian
coefficients at base $1/4$.
\end{corollary}

\section{The endpoint correction is a Thue--Morse logarithm}
\label{sec:endpoint-TM}

\subsection{Exact identification}

\begin{theorem}[Mahler identification of the Laplace tail]\label{thm:tail-TM}
\statusN\ For every $s>0$,
\begin{equation}
 \boxed{
 \e^{T(s)}=\Theta(\e^{-s})
 =\sum_{n\ge0}\varepsilon_n\e^{-ns},
 \qquad
 \mathcal E(s)=-\log\Theta(\e^{-s}).}                      \label{eq:tail-TM}
\end{equation}
Consequently the exact periodic decomposition \eqref{eq:Lambda-periodic} becomes
\begin{equation}
 \boxed{
 \Lambda(s)=-\frac{(\log s)^2}{2L}+\frac12\log s
 +P(\log_2s)-\log\Theta(\e^{-s}).}                         \label{eq:Lambda-TM}
\end{equation}
\end{theorem}

\begin{proof}
Substitute $q=\e^{-s}$ in the Thue--Morse product \eqref{eq:Theta-product}:
\[
 \Theta(\e^{-s})=\prod_{m\ge0}(1-\e^{-2^ms})=\e^{T(s)}.
\]
All factors are positive, so the real logarithm is unambiguous.  The remaining formulas are
immediate from \eqref{eq:E-tail} and \eqref{eq:Lambda-periodic}.
\end{proof}

\subsection{The $2$-adic Lambert series}
Define
\begin{equation}
 a(n)=\sum_{j=0}^{\vtwo(n)}2^j=2^{\vtwo(n)+1}-1.           \label{eq:a-n}
\end{equation}

\begin{theorem}[Lambert expansion and derivative hierarchy]\label{thm:Lambert-tail}
\statusN\ For $s>0$,
\begin{align}
 \mathcal E(s)&=\sum_{n\ge1}\frac{a(n)}n\e^{-ns},         \label{eq:E-Lambert}\\
 T^{(k)}(s)&=(-1)^{k+1}\sum_{n\ge1}a(n)n^{k-1}\e^{-ns}
 \qquad(k\ge1).                                           \label{eq:T-derivatives}
\end{align}
In particular, with
\begin{equation}
 R(s)=sT'(s),                                               \label{eq:R-def}
\end{equation}
one has
\begin{equation}
 \boxed{
 R(s)=\sum_{m\ge0}\frac{2^ms}{\e^{2^ms}-1}
 =s\sum_{n\ge1}a(n)\e^{-ns}.}                             \label{eq:R-Lambert}
\end{equation}
The beginning of the endpoint correction is
\begin{equation}
 \mathcal E(s)=\e^{-s}+\frac32\e^{-2s}+\frac13\e^{-3s}
 +\frac74\e^{-4s}+\frac15\e^{-5s}+\cdots.               \label{eq:E-first}
\end{equation}
\end{theorem}

\begin{proof}
Expand every logarithm in \eqref{eq:E-tail}:
\[
 \mathcal E(s)=\sum_{m\ge0}\sum_{k\ge1}\frac{\e^{-k2^ms}}k.
\]
For a fixed $n$, the representations $n=k2^m$ are indexed by
$m=0,\ldots,\vtwo(n)$, and $1/k=2^m/n$.  Their sum is $a(n)/n$.
Termwise differentiation is justified by absolute convergence on every half-line
$s\ge s_0>0$.
\end{proof}

\subsection{The spectral--endpoint dictionary}

\begin{theorem}[Endpoint coefficients from Fourier zero multiplicities]
\label{thm:spectral-endpoint}
\statusN\ Let $\mu(n)=\ord_{z=n}\Phi$.  Then
\begin{equation}
 a(n)=2^{\mu(n)}-1,                                        \label{eq:a-mu}
\end{equation}
and hence
\begin{equation}
 \boxed{
 -\log\Theta(\e^{-s})
 =\sum_{n\ge1}\frac{2^{\ord_{z=n}\Phi}-1}{n}\e^{-ns}.}   \label{eq:spectral-endpoint}
\end{equation}
Equivalently, the weighted zero-multiplicity Dirichlet series is
\begin{equation}
 \boxed{
 \mathfrak Z_\Phi(w):=
 \sum_{n\ge1}\frac{2^{\ord_{z=n}\Phi}-1}{n^w}
 =\frac{\zeta(w)}{1-2^{1-w}},\qquad \Re w>1.}             \label{eq:weighted-zero-zeta}
\end{equation}
\end{theorem}

\begin{proof}
By \eqref{eq:zero-mult}, $\mu(n)=\vtwo(n)+1$, so \eqref{eq:a-mu} is
\eqref{eq:a-n}.  For the Dirichlet series, write
$a(n)=\sum_{j\ge0}2^j\mathbf1_{2^j\mid n}$ and use Tonelli:
\[
 \sum_{n\ge1}a(n)n^{-w}
 =\zeta(w)\sum_{j\ge0}2^{j(1-w)}
 =\frac{\zeta(w)}{1-2^{1-w}}.
\]
\end{proof}

\begin{boxedremark}
\textbf{Why this is a genuine bridge.}
The left side of \eqref{eq:spectral-endpoint} belongs to the small-$x$ Fabius endpoint
problem; the multiplicity on the right belongs to the zero divisor of the Rvachev Fourier
transform.  The equality does not merely compare asymptotic scales: it identifies every
individual exponential coefficient with a local spectral multiplicity.
\end{boxedremark}

\begin{center}
\begin{tabular}{@{}rrrr@{}}
\toprule
$n$ & $\vtwo(n)$ & $\ord_{z=n}\Phi$ & coefficient of $\e^{-ns}$ in $\mathcal E(s)$\\
\midrule
1 & 0 & 1 & $1$\\
2 & 1 & 2 & $3/2$\\
3 & 0 & 1 & $1/3$\\
4 & 2 & 3 & $7/4$\\
6 & 1 & 2 & $1/2$\\
8 & 3 & 4 & $15/8$\\
12 & 2 & 3 & $7/12$\\
16 & 4 & 5 & $31/16$\\
\bottomrule
\end{tabular}
\end{center}

\section{Mellin transform and the exact periodic phase}
\label{sec:Mellin}

\subsection{Mellin transform}

\begin{theorem}[Positive-strip realization of the Gamma--zeta kernel]
\label{thm:Mellin-tail}
\statusN\ For $\Re w>0$,
\begin{equation}
 \boxed{
 \Mellin\mathcal E(w):=\int_0^\infty\mathcal E(s)s^{w-1}\dd s
 =\frac{\Gamma(w)\zeta(w+1)}{1-2^{-w}}.}                   \label{eq:Mellin-tail}
\end{equation}
Equivalently, the Mellin transform of $R(s)=sT'(s)=-s\mathcal E'(s)$ is
\begin{equation}
 \int_0^\infty R(s)s^{w-1}\dd s
 =\frac{\Gamma(w+1)\zeta(w+1)}{1-2^{-w}}.                 \label{eq:Mellin-R}
\end{equation}
\statusE\ The corpus obtains the same meromorphic right-hand side as the Mellin transform
of $\Lambda$ on the complementary fundamental strip $-1<\Re w<0$.  Thus
\eqref{eq:Mellin-tail} is not a new special-function expression; its corpus-relative
novelty is that the Thue--Morse endpoint tail and the weighted Fourier zero divisor realize
that established kernel on $\Re w>0$.
\end{theorem}

\begin{proof}
Put $\sigma=\Re w>0$.  Absolute convergence permits Fubini, because
\[
 \sum_{n\ge1}\frac{a(n)}n
 \int_0^\infty\e^{-ns}s^{\sigma-1}\dd s
 =\Gamma(\sigma)\sum_{n\ge1}\frac{a(n)}{n^{\sigma+1}}<\infty.
\]
Therefore \eqref{eq:E-Lambert} gives
\[
 \int_0^\infty\mathcal E(s)s^{w-1}\dd s
 =\Gamma(w)\sum_{n\ge1}\frac{a(n)}{n^{w+1}}.
\]
Use \eqref{eq:weighted-zero-zeta} with $w+1$ in place of $w$.  Integration by parts gives
\eqref{eq:Mellin-R}; the boundary terms vanish for $\Re w>0$.
\end{proof}

\subsection{Residues and the Fourier coefficients of $P$}
Put
\begin{equation}
 \chi_k=\frac{2\pi\ii k}{L},\qquad k\in\Z.                 \label{eq:chi}
\end{equation}
We use the Fourier convention
$P(u)=\sum_{k\in\Z}\widehat P(k)\e^{2\pi\ii ku}$.

\begin{theorem}[Mellin recovery of the full periodic spectrum]\label{thm:P-Fourier}
\statusE\ The periodic function in \eqref{eq:Lambda-periodic} has
\begin{equation}
 \boxed{
 \widehat P(k)=-\frac{\Gamma(-\chi_k)\zeta(1-\chi_k)}L,
 \qquad k\ne0,}                                           \label{eq:P-nonzero}
\end{equation}
which agrees with the nonzero-mode formula already established in the corpus.  Its mean is
\begin{equation}
 \boxed{
 \widehat P(0)
 =-\frac L{12}
 -\frac1L\left(\frac{\pi^2}{12}-\gamma_1-\frac{\gamma^2}{2}\right),} \label{eq:P-mean}
\end{equation}
where $\gamma$ is Euler's constant and $\gamma_1$ is the first Stieltjes constant in the
convention
$\zeta(1+w)=w^{-1}+\gamma-\gamma_1w+\bigO(w^2)$.
Numerically,
\begin{equation}
 \widehat P(0)=-1.1090453654341866821907455547\ldots.       \label{eq:P-mean-num}
\end{equation}
These coefficient formulas, including the mean, are established in the corpus (the mean is
formalized there).  \statusN\ The corpus-relative novelty is the unified derivation below:
both the oscillatory modes and the zero mode arise from the same Thue--Morse endpoint
correction and the same meromorphic Mellin kernel.
\end{theorem}

\begin{proof}
Mellin inversion writes $\mathcal E(s)$ as a vertical integral of
\[
 M(w)s^{-w},\qquad M(w)=\frac{\Gamma(w)\zeta(w+1)}{1-2^{-w}}.
\]
The exponential decay of $\Gamma$ on vertical lines justifies shifting the contour across
$\Re w=0$ and summing the resulting residues absolutely.  Besides the pole at $w=0$, the
denominator has simple poles at $w=\chi_k$.
The pole at $w=-\chi_k$ contributes
\[
 \frac{\Gamma(-\chi_k)\zeta(1-\chi_k)}L\,s^{\chi_k}
 =\frac{\Gamma(-\chi_k)\zeta(1-\chi_k)}L\,\e^{2\pi\ii k\log_2s}.
\]
On the other hand, \eqref{eq:Lambda-periodic} and $\Lambda(s)=\bigO(s)$ as $s\downarrow0$
give
\begin{equation}
 \mathcal E(s)=\frac{(\log s)^2}{2L}-\frac12\log s
 -P(\log_2s)+\bigO(s).                                    \label{eq:E-small}
\end{equation}
Comparison proves \eqref{eq:P-nonzero}.

For the zero mode, expand at $w=0$:
\begin{align*}
 \Gamma(w)\zeta(1+w)
 &=\frac1{w^2}+c_0+\bigO(w),
 &c_0&=\frac{\pi^2}{12}-\gamma_1-\frac{\gamma^2}{2},\\
 \frac1{1-2^{-w}}
 &=\frac1{Lw}+\frac12+\frac{Lw}{12}+\bigO(w^3).
\end{align*}
Hence
\[
 M(w)=\frac1{Lw^3}+\frac1{2w^2}
 +\left(\frac L{12}+\frac{c_0}L\right)\frac1w+\bigO(1).
\]
The residue of $M(w)s^{-w}$ at zero is
$(\log s)^2/(2L)-(\log s)/2+L/12+c_0/L$.  Comparing with
\eqref{eq:E-small} shows that this constant is $-\widehat P(0)$.
\end{proof}

\begin{remark}[Why the oscillation is easy to miss]
The first nonzero mode is tiny:
\[
 |\widehat P(1)|=1.0627116252\ldots\times10^{-6},\qquad
 |\widehat P(2)|=6.6927864\ldots\times10^{-13}.
\]
Differentiation magnifies the first mode, but it is still small.  This explains why a
non-periodic numerical fit can look excellent while being wrong at the sharp error scale.
\end{remark}

\subsection{The periodic correction to the saddle mean}
Define
\begin{equation}
 \Pi(u)=-\frac12-\frac{P'(u)}L.                             \label{eq:Pi}
\end{equation}
Then $\widehat\Pi(0)=-1/2$, while for $k\ne0$,
\begin{equation}
 \boxed{
 \widehat\Pi(k)=-\frac{\Gamma(1-\chi_k)\zeta(1-\chi_k)}L.} \label{eq:Pi-Fourier}
\end{equation}
The first harmonic has magnitude $9.6331836245\ldots\times10^{-6}$.

\section{Tilted cumulants and exact periodically perturbed Lambert saddles}
\label{sec:saddle}

\subsection{Exact tilted mean}
Under the exponentially tilted law
\[
 \frac{\e^{-sX}}{B(s)}\dd\Prob,
\]
write $\E_s$ and $\Var_s$ for expectation and variance.  Define
\begin{equation}
 A(s)=-s\Lambda'(s)=s\E_sX.                                \label{eq:A}
\end{equation}

\begin{theorem}[Mean decomposition and dyadic recurrence]\label{thm:A-decomp}
\statusN\ For every $s>0$,
\begin{equation}
 \boxed{A(s)=\log_2s+\Pi(\log_2s)+R(s).}                   \label{eq:A-decomp}
\end{equation}
It also satisfies the exact recurrence
\begin{equation}
 \boxed{A(2s)=A(s)+1-\frac{s}{\e^s-1}.}                    \label{eq:A-recurrence}
\end{equation}
\end{theorem}

\begin{proof}
Apply $-s\frac\dd{\dd s}$ to \eqref{eq:Lambda-periodic}; this gives
\eqref{eq:A-decomp}.  Differentiate \eqref{eq:Lambda-dilation} and multiply by $-s$ to
obtain \eqref{eq:A-recurrence}.
\end{proof}

\subsection{All tilted cumulants}
Let
\begin{equation}
 \kappa_n(s)=(-1)^n\Lambda^{(n)}(s),\qquad
 K_n(s)=s^n\kappa_n(s),\qquad D=s\frac\dd{\dd s}.           \label{eq:tilted-cumulants}
\end{equation}
Thus $\kappa_n(s)$ is the $n$th cumulant of $X$ under the tilted law.

\begin{theorem}[Falling-Euler hierarchy]\label{thm:falling-Euler}
\statusN\ For $n\ge1$,
\begin{equation}
 \boxed{
 K_n(s)=(-1)^n\falling Dn
 \left[-\frac L2u(u-1)+P(u)-T(s)\right],
 \qquad u=\log_2s,}                                       \label{eq:falling-Euler}
\end{equation}
where $\falling Dn=D(D-1)\cdots(D-n+1)$.  If
$h(s)=\log(1-\e^{-s})-\log s$, then
\begin{equation}
 \boxed{K_n(2s)=K_n(s)+(-1)^ns^nh^{(n)}(s).}               \label{eq:Kn-dilation}
\end{equation}
In particular,
\begin{equation}
 s^2\Var_s(X)
 =A(s)-DA(s)
 =u+\Pi(u)+R(s)-\frac{1+\Pi'(u)}L-sR'(s)>0,                \label{eq:variance-exact}
\end{equation}
and
\begin{equation}
 K_2(2s)=K_2(s)+1-\frac{s^2\e^s}{(\e^s-1)^2}.             \label{eq:K2-recurrence}
\end{equation}
\end{theorem}

\begin{proof}
The operator identity $s^n\frac{\dd^n}{\dd s^n}=\falling Dn$ proves
\eqref{eq:falling-Euler}.  Differentiate \eqref{eq:Lambda-dilation} $n$ times to obtain
\eqref{eq:Kn-dilation}.  Since $K_2=s^2\Lambda''=A-DA$, substitution of
\eqref{eq:A-decomp} gives \eqref{eq:variance-exact}.  Strict positivity follows because
$X$ is nondegenerate under every finite tilt.  The explicit formula
\eqref{eq:K2-recurrence} follows by differentiating $h$ twice.
\end{proof}

\subsection{Exact Lambert fixed point and branch threshold}
The natural saddle for the inverse Laplace problem at $x\in(0,1/2)$ is the unique solution
of
\begin{equation}
 x=-\Lambda'(s)=\E_sX.                                    \label{eq:saddle-s}
\end{equation}
Write $s=2^u$ and
\begin{equation}
 C(u)=\Pi(u)+R(2^u)=A(2^u)-u.                              \label{eq:C-u}
\end{equation}
Then the exact saddle equation is
\begin{equation}
 x2^u=u+C(u).                                               \label{eq:saddle-u}
\end{equation}

\begin{theorem}[Periodically perturbed Lambert maps]\label{thm:periodic-W}
\statusN\ Equation \eqref{eq:saddle-u} has one solution $u=u_*(x)$ for every
$0<x<1/2$.  There is a unique $s_c>0$ such that
\begin{equation}
 A(s_c)=\frac1L,
 \qquad
 x_c:=\frac{A(s_c)}{s_c}=\frac1{Ls_c},                    \label{eq:xc-def}
\end{equation}
Numerically,
\begin{equation}
 s_c=3.5627054762160828049\ldots,
 \qquad x_c=0.4049436728688661444\ldots.                  \label{eq:xc-num}
\end{equation}
For the endpoint range $0<x<x_c$, the saddle is equivalently the fixed point of the
lower-branch map
\begin{equation}
 \boxed{
 u=\mathfrak L_x^{(-1)}(u):=-\frac1L
 W_{-1}\!\left(-Lx\,2^{-C(u)}\right)-C(u).}                \label{eq:W-map}
\end{equation}
For $x_c<x<1/2$, the same formula with $W_0$ in place of $W_{-1}$ gives the saddle; at
$x=x_c$ the two branches meet.  At a fixed point of either branch,
\begin{equation}
 \boxed{
 \bigl(\mathfrak L_x^{(b)}\bigr)'(u)
 =\frac{C'(u)}{L(u+C(u))-1},\qquad b\in\{0,-1\}.}          \label{eq:W-derivative}
\end{equation}
Because $\Pi'$ is bounded and $R(2^u)$ is exponentially small as $u\to\infty$, the
lower-branch map is a local contraction at $u_*(x)$ for all sufficiently small $x$.
\end{theorem}

\begin{proof}
The function $s\mapsto\E_sX$ is strictly decreasing because its derivative is
$-\Var_s(X)<0$.  Its limits are $1/2$ at $s\downarrow0$ and $0$ at $s\to\infty$, proving
existence and uniqueness of the saddle.  Also
\[
 A(s)=\sum_{j\ge1}g(s/2^j),\qquad
 g(t)=1-\frac{t}{\e^t-1},
\]
and
$g'(t)=\bigl(\e^t(t-1)+1\bigr)/(\e^t-1)^2>0$ for $t>0$.
Thus $A$ increases continuously from $0$ to infinity (for instance, the first
$\lfloor\log_2s\rfloor$ summands are uniformly positive once $s$ is large), so
\eqref{eq:xc-def} has one solution.  Since $x=A(s)/s=\E_sX$ decreases with $s$, one has
$A(s)>1/L$ exactly when $x<x_c$.

Put $y=u+C(u)=A(2^u)$.  Equation \eqref{eq:saddle-u} becomes
$y\e^{-Ly}=x\e^{-LC(u)}$.  Multiplication by $-L$ selects $W_{-1}$ when
$y>1/L$, $W_0$ when $0<y<1/L$, and the common branch value at $y=1/L$.  This proves the
branch statements.  Implicit differentiation of
$y\e^{-Ly}=x\e^{-LC(u)}$, followed by subtraction of $C'(u)$, gives
\eqref{eq:W-derivative}.  In the endpoint regime the denominator tends to infinity,
whereas $C'$ stays bounded.
\end{proof}

\begin{corollary}[First correction to the unperturbed Lambert phase]
Let
\begin{equation}
 u_0(x)=-\frac1L W_{-1}(-Lx),                              \label{eq:u0}
\end{equation}
so that $x2^{u_0}=u_0$.  Then, as $x\downarrow0$,
\begin{equation}
 u_*(x)=u_0(x)+\frac{C(u_0(x))}{Lu_0(x)-1}
 +\bigO_C\!\left(u_0(x)^{-2}\right).                       \label{eq:first-W-correction}
\end{equation}
A single evaluation of the exact map \eqref{eq:W-map} is typically much more accurate than
the linear correction.
\end{corollary}

\begin{proofidea}
Apply the mean-value theorem to
$f(u)=x2^u-u-C(u)$ between $u_0$ and $u_*$.  Since
$f(u_0)=-C(u_0)$ and $f'(u)=Lx2^u-1-C'(u)$, the displacement is of order $1/u_0$.
A second expansion gives \eqref{eq:first-W-correction}; the implied constant depends only on uniform bounds for $C$ and its first two derivatives.
\end{proofidea}

\subsection{Numerical check of the exact fixed point}
The following values were computed at 70 decimal digits by summing the exact tilted mean
\[
 A(s)=\sum_{j\ge1}\left(1-\frac{s/2^j}{\e^{s/2^j}-1}\right)
\]
and solving $x2^u=A(2^u)$.
\begin{center}
\small
\begin{tabular}{@{}lrrrr@{}}
\toprule
$x$ & $u_0$ & $u_*$ & linear-correction error & one-step map error\\
\midrule
$10^{-3}$
& $13.7468091666470$ & $13.6868068298588$
& $1.3742\times10^{-3}$ & $7.4062\times10^{-7}$\\
$10^{-6}$
& $24.5491709694670$ & $24.5175886975057$
& $3.6440\times10^{-4}$ & $2.0539\times10^{-7}$\\
$10^{-12}$
& $45.3666989202317$ & $45.3501796089457$
& $9.7317\times10^{-5}$ & $6.5109\times10^{-9}$\\
\bottomrule
\end{tabular}
\end{center}
The leading displacement is negative because $C(u)=-1/2+$ a tiny periodic oscillation plus
an exponentially small term.

\section{All-orders sinc tails and centered Rvachev moments}
\label{sec:tail-moments}

\subsection{Exact base-$1/4$ tail law}
The finite product $\Phi_m$ from \eqref{eq:Phi-m-Sm} satisfies
\begin{equation}
 \frac{\Phi(z)}{\Phi_m(z)}=\Phi(z/2^m).                    \label{eq:tail-shell}
\end{equation}

\begin{theorem}[All-orders tail logarithm]\label{thm:tail-log}
\statusN\ If $|z|<2^m$ and
\begin{equation}
 Q_m(z)=\frac{z^2}{4^{m-1}},                               \label{eq:Qm}
\end{equation}
then
\begin{equation}
 \boxed{
 \log\frac{\Phi(z)}{\Phi_m(z)}
 =-\sum_{r\ge1}\frac{\zeta(2r)}{r(4^r-1)}Q_m(z)^r.}       \label{eq:tail-log}
\end{equation}
Writing
\begin{equation}
 \frac{\Phi(z)}{\Phi_m(z)}=\sum_{N\ge0}a_NQ_m(z)^N,
 \qquad a_0=1,                                             \label{eq:aN-def}
\end{equation}
the coefficients satisfy
\begin{equation}
 \boxed{
 Na_N=-\sum_{r=1}^{N}\frac{\zeta(2r)}{4^r-1}a_{N-r}.}     \label{eq:aN-rec}
\end{equation}
The first terms are
\begin{align}
 1-\frac{\pi^2}{18}Q
 +\frac{19\pi^4}{16200}Q^2
 -\frac{583\pi^6}{42865200}Q^3
 +\frac{132809\pi^8}{1311675120000}Q^4+\cdots.            \label{eq:tail-first}
\end{align}
\end{theorem}

\begin{proof}
Use \eqref{eq:q-log-full} with $q=1/4$ and replace $z$ by $z/2^m$.
The coefficient recurrence follows by logarithmic differentiation of
\eqref{eq:aN-def}.
\end{proof}

\subsection{Centered cumulants and a Bernoulli--Mersenne recurrence}
Let
\begin{equation}
 c_N=\int_{-1}^{1}x^{2N}\Up(x)\dd x=\E Y^{2N}.             \label{eq:cN}
\end{equation}

\begin{theorem}[Tail coordinates for centered Rvachev moments]\label{thm:centered-cumulants}
\statusE\ The corpus's centered law has cumulants
\begin{equation}
 \boxed{
 \kappa_{2r}(Y)=\frac{2^{2r-1}B_{2r}}{r(4^r-1)},
 \qquad \kappa_{2r+1}(Y)=0,}                               \label{eq:centered-cumulants}
\end{equation}
and its centered moments $c_N$ are already used in the exact dyadic theory.
\statusN\ In the spectral coordinates of \cref{thm:tail-log}, those moments satisfy the
alternative Bernoulli--Mersenne recurrence
\begin{equation}
 \boxed{
 c_N=\frac1N\sum_{r=1}^{N}\binom{2N}{2r}
 \frac{2^{2r-1}B_{2r}}{4^r-1}\,c_{N-r},}                 \label{eq:centered-moment-rec}
\end{equation}
and the tail coefficient is exactly
\begin{equation}
 \boxed{a_N=\frac{(-1)^Nc_N\pi^{2N}}{(2N)!}.}             \label{eq:aN-cN}
\end{equation}
\end{theorem}

\begin{proof}
Specialize \cref{thm:q-cumulants} to $q=1/4$.  The recurrence is the corresponding
specialization of \eqref{eq:q-moment-rec}.  Finally
\[
 \Phi(z/2^m)=\sum_{N\ge0}\frac{(-1)^Nc_N}{(2N)!}
 \left(\frac{\pi z}{2^{m-1}}\right)^{2N},
\]
and comparison with \eqref{eq:aN-def} gives \eqref{eq:aN-cN}.
\end{proof}

The first moments are
\begin{equation}
 c_0=1,
 \quad c_1=\frac19,
 \quad c_2=\frac{19}{675},
 \quad c_3=\frac{583}{59535},
 \quad c_4=\frac{132809}{32531625},
 \quad c_5=\frac{46840699}{24405225075}.                   \label{eq:c-first}
\end{equation}
The variance is $1/9$, and the fourth cumulant is $-2/225$.

\begin{remark}
The corpus already defines the centered moments, gives their hyperbolic-sine generating
product, tabulates the values in \eqref{eq:c-first}, and proves a different Bernoulli
recurrence.  The new contribution here is narrower: the general-$q$ derivation, the
alternative cumulant-form recurrence \eqref{eq:centered-moment-rec}, the exact identification
of these established moments with the base-$1/4$ sinc-tail coefficients, and the certified
Thue--Morse acceleration below.
\end{remark}

\subsection{Certified Thue--Morse acceleration}
Combining the finite triple factorization with the moment tail gives an exact hybrid
representation.

\begin{theorem}[Thue--Morse spectral reconstruction]\label{thm:TM-reconstruction}
\statusN\ For every $m\ge1$ and $z\in\C$,
\begin{align}
 \Phi(z)
 &=\frac{2^{m(m-3)/2}\e^{2\pi\ii z(1-2^{-m})}}{(\ii\pi z)^m}
 S_m(z)                                                     \notag\\
 &\quad\times
 \sum_{N=0}^{\infty}\frac{(-1)^Nc_N}{(2N)!}
 \left(\frac{\pi z}{2^{m-1}}\right)^{2N}.                 \label{eq:TM-reconstruction}
\end{align}
For real $z$, truncation after $N=M$ has the certified error
\begin{align}
 &\left|\Phi(z)-\Phi_m(z)
 \sum_{N=0}^{M}\frac{(-1)^Nc_N}{(2N)!}
 \left(\frac{\pi z}{2^{m-1}}\right)^{2N}\right|           \notag\\
 &\hspace{35mm}\le
 \frac{(\pi|z|/2^{m-1})^{2M+2}}{(2M+2)!}.                 \label{eq:TM-error}
\end{align}
\end{theorem}

\begin{proof}
Multiply \eqref{eq:triple-factor} by the moment series for
$\Phi(z/2^m)$.  For real $z$, interpret the tail as
$\E\cos((\pi z/2^{m-1})Y)$ and use the Lagrange remainder for cosine together with
$|Y|\le1$ and $|\Phi_m(z)|\le1$.
\end{proof}

\begin{example}[Numerical acceleration]
At $z=10$ and $m=6$, the exact tail is
$\Phi(10/64)=0.94753140589767634644\ldots$.  The errors after retaining $M$ moment
corrections, together with the rigorous bound \eqref{eq:TM-error}, are
\begin{center}
\begin{tabular}{@{}rrr@{}}
\toprule
$M$ & actual tail error & certified upper bound\\
\midrule
0 & $5.24686\times10^{-2}$ & $4.81914\times10^{-1}$\\
1 & $1.07744\times10^{-3}$ & $3.87069\times10^{-2}$\\
2 & $1.20907\times10^{-5}$ & $1.24356\times10^{-3}$\\
3 & $8.69394\times10^{-8}$ & $2.14032\times10^{-5}$\\
\bottomrule
\end{tabular}
\end{center}
Thus a finite signed Thue--Morse sum, corrected by only a few rational moments, gives a
rapid and rigorously controlled evaluation of the infinite sinc product.
\end{example}

\section{How the new identities reorganize the existing theory}

\subsection{A commutative diagram of transforms}
The central relationships can be read as the following chain:
\[
\begin{array}{ccccc}
 \Theta(q)=\prod(1-q^{2^j})
 &\xrightarrow{\ q=\e^{-s}\ }&
 \mathcal E(s)=-\log\Theta(\e^{-s})
 &\xrightarrow{\ \Mellin\ }&
 \dfrac{\Gamma(w)\zeta(w+1)}{1-2^{-w}}\\[1.2em]
 &&\Big\updownarrow&&\Big\downarrow\text{ residues}\\[-0.2em]
 \Phi(z)=\prod\sinc(\pi z/2^j)
 &\xrightarrow{\ \text{zero divisor}\ }&
 \mu(n)=\vtwo(n)+1
 &\xrightarrow{\ a(n)=2^{\mu(n)}-1\ }&
 P(\log_2s).
\end{array}
\]
The periodic endpoint function is therefore not an unrelated analytic decoration.  It is
the residue shadow of a Thue--Morse logarithm whose coefficients can be read from the
Rvachev spectral zero divisor.

\subsection{Two compatible $q$-geometries}
The corpus's exact dyadic arithmetic uses base $q=1/2$ because each spatial refinement
changes the dyadic denominator by one power of two.  The Fourier product uses base
$q=1/4$ because its Taylor and Euler products are even in $z$, so the same spatial
dilation becomes a squared spectral dilation.  This explains why
\begin{itemize}
\item $(1/2;1/2)_n$ and $\qbinom nk{1/2}$ arise in exact Fabius values;
\item $(z^2/\ell^2;1/4)_\infty$, $4^r-1$, and
$\qbinom mj{1/4}$ arise in the sinc product and centered moments.
\end{itemize}
A future theory should treat these not as competing normalizations but as the arithmetic
and spectral faces of the same binary refinement.

\subsection{Endpoint asymptotics from spectral data}
The exact identity \eqref{eq:spectral-endpoint} gives a new route toward the endpoint
expansion.  Starting only with the Fourier zero multiplicities,
\begin{enumerate}
\item form $a(n)=2^{\mu(n)}-1$;
\item build $\mathcal E(s)=\sum a(n)\e^{-ns}/n$;
\item Mellin transform to obtain \eqref{eq:Mellin-tail};
\item read off the quadratic logarithm and all periodic modes;
\item insert them into the exact saddle map \eqref{eq:W-map}.
\end{enumerate}
This does not yet derive every coefficient of the all-orders Fabius endpoint expansion,
but it supplies the entire oscillatory datum and the exact periodic saddle displacement
from the spectral zero divisor alone.

\subsection{Relation to the Fourier-decay audits}
The arithmetic product \eqref{eq:arithmetic-product} is not by itself a pointwise asymptotic
for $|\Phi(x)|$: near integers, zero multiplicities dominate, while extremal rays avoid
zeros in highly structured ways.  It nevertheless offers three tools for the audited
decay problem:
\begin{enumerate}
\item exact zero multiplicities and cumulative zero counting;
\item a $q$-family in which the binary ratio can be varied continuously;
\item certified finite-shell corrections suitable for rigorous numerics and Lean
formalization.
\end{enumerate}
The separation of $\kappa_\infty,\kappa_2,\kappa_1,\kappa_0$ remains essential.

\section{Further frontier problems}
\label{sec:frontiers}

\subsection{High-priority formalization targets}
The following statements appear especially suitable for Lean because their proofs reduce to
finite products, absolutely convergent series, or elementary operator identities.
\begin{enumerate}
\item \textbf{Arithmetic canonical product:} formalize
$\Phi(z)=\prod_{n\ge1}(1-z^2/n^2)^{\vtwo(n)+1}$ and the integer zero multiplicity.
\item \textbf{Finite triple factorization:} combine the finite sinc product with the
formal Thue--Morse polynomial and Prouhet cancellation.
\item \textbf{Endpoint Mahler identity:} prove
$\mathcal E(s)=-\log\Theta(\e^{-s})$ and the $2$-adic Lambert expansion.
\item \textbf{Spectral--endpoint coefficient theorem:} identify the Lambert coefficient
with $2^{\ord_n\Phi}-1$.
\item \textbf{Centered moment recurrence:} derive
\eqref{eq:centered-moment-rec} from the already formal sinc product.
\item \textbf{Periodic Lambert map:} package the exact saddle equation and derivative
\eqref{eq:W-derivative}, using strict tilted variance for uniqueness.
\end{enumerate}
The Mellin residue computation and Stieltjes-constant mean are analytically heavier but can
be isolated after the elementary core is formalized.

\subsection{A general-$q$ decay spectrum}
\begin{conjecture}[Smooth dependence of decay exponents]\label{conj:q-spectrum}
\statusC\ For the geometric product $\Phi_q$, the analogues of the extremal, $L^2$,
Ces\`aro, and typical power exponents exist for every $0<q<1$ and vary continuously, perhaps
real-analytically away from resonance parameters at which zero lattices overlap.
\end{conjecture}
The universal quadratic coefficient should be
$1/(2|\log\sqrt q|)=1/|\log q|$, as the elementary upper bound in
\cref{thm:q-density} already suggests.  The lower-order exponents should arise from transfer
operators depending on $q$.

\subsection{Bridging the bases $1/2$ and $1/4$}
\begin{conjecture}[Arithmetic--spectral $q$-transform]
\statusC\ There is a direct coefficient transform carrying the base-$1/2$ Gaussian-binomial
formulas for exact dyadic values of $F$ to the base-$1/4$ Bernoulli--Mersenne moments of
$\Up$, without passing through the full probability model.
\end{conjecture}
A likely starting point is the identity $Y=2X-1$, followed by a binomial transform of the
half moments, but a satisfactory formula should preserve the $q$-Pochhammer structure
rather than expand it away.

\subsection{A spectral derivation of the full endpoint expansion}
\begin{conjecture}[Zero-divisor endpoint reconstruction]
\statusC\ The complete all-orders endpoint asymptotic of $F(x)$, including every periodic
coefficient, can be reconstructed from the weighted zero zeta function
$\mathfrak Z_\Phi$ together with the local normalization $\Phi(0)=1$.
\end{conjecture}
The report reconstructs the quadratic logarithm and rederives the entire periodic function
$P$, including its established mean, from the weighted zero data; it also derives the exact
saddle displacement.  The remaining challenge is to derive the full
Gaussian saddle-jet algebra directly from spectral data.

\subsection{Arithmetic nature of the periodic mean}
Formula \eqref{eq:P-mean} introduces the combination
\[
 \frac{\pi^2}{12}-\gamma_1-\frac{\gamma^2}{2}.
\]
Its arithmetic nature is unknown.  Even without transcendence information, the exact
formula is useful: it separates the mean from the exponentially tiny nonzero modes and
eliminates one numerical constant from the endpoint theory.

\subsection{Weighted zero statistics}
The ordinary multiplicity and endpoint-weighted multiplicity have Dirichlet series
\[
 \sum\frac{\mu(n)}{n^w}=\frac{\zeta(w)}{1-2^{-w}},
 \qquad
 \sum\frac{2^{\mu(n)}-1}{n^w}=\frac{\zeta(w)}{1-2^{1-w}}.
\]
The second has a double singularity at $w=1$ after multiplication by the pole of $\zeta$,
which is the arithmetic source of the quadratic logarithm in the endpoint correction.
A Tauberian development of its summatory function, with explicit binary periodic
fluctuation, should give an elementary real-variable route to the Mellin result.

\section{Conclusions}
The inspected corpus already contains an unusually complete and carefully status-separated
theory of the Fabius and Rvachev functions.  The main contribution of this report is not a
new isolated formula but a new exact dictionary:
\begin{enumerate}
\item the infinite sinc product is a base-$1/4$ $q$-Pochhammer product;
\item its integer zero multiplicity is $\vtwo(n)+1$;
\item the Thue--Morse endpoint logarithm weights that multiplicity by $2^\mu-1$;
\item the resulting Dirichlet series has Mellin transform
$\Gamma(w)\zeta(w+1)/(1-2^{-w})$;
\item the residues reproduce the endpoint's full established log-periodic spectrum,
including its formalized mean, through a new Thue--Morse/Mellin route;
\item the same periodic function enters an exact lower-$W_{-1}$ fixed point for the saddle;
\item finite Thue--Morse sums, corrected by centered Bernoulli--Mersenne moments, give
certified approximations to the full Fourier product.
\end{enumerate}
This unifies the Thue--Morse, $q$-binomial, sinc-product, endpoint, and Lambert-$W$ themes
without erasing their different normalizations.  It also leaves a concrete formalization
program and several testable frontier questions.

\appendix
\section{Supplementary derivations}

\subsection{The coefficient formula by partitions}
From the cumulant generating function, the centered moments also admit the explicit
partition sum
\begin{equation}
 c_N=(2N)!\!\sum_{\substack{m_1,\ldots,m_N\ge0\\
                    \sum_{r=1}^Nrm_r=N}}
 \prod_{r=1}^N\frac1{m_r!}
 \left[
  \frac{2^{2r-1}B_{2r}}{r(4^r-1)(2r)!}
 \right]^{m_r}.                                            \label{eq:partition-moments}
\end{equation}
This makes rationality immediate and displays the product of Mersenne-type factors
$4^r-1$.  Since
\[
 \prod_{r=1}^N(4^r-1)=4^{N(N+1)/2}(1/4;1/4)_N,
\]
the denominator geometry is again governed by a base-$1/4$ $q$-Pochhammer symbol.

\paragraph{Post-snapshot Lean status.}
The universal algebraic shape of \eqref{eq:partition-moments} is now
formalized in \texttt{ExponentialPartition.lean}: with
\[
 E_r=\frac{2^{2r-1}B_{2r}}{r(4^r-1)(2r)!},
\]
\texttt{partitionExpSum\_eq\_sum\_div} identifies
\(\texttt{partitionExpSum}\ E\ N\) with the displayed partition sum
before multiplication by \((2N)!\).
\texttt{ExponentialBell.lean} further identifies that coefficient
function with \texttt{SaddleExpansion.expCoeff}.  These are generic
coefficient identities only: the cumulant formula, its identification
with \(c_N\), the specialized Bernoulli--Mersenne recurrence, and a
standalone complete-Bell-polynomial interface remain separate formal
obligations.  The \statusN\ classifications above therefore remain
unchanged.

\subsection{A direct proof of the first-harmonic scale}
For $k\ne0$, Stirling's formula on vertical lines gives
\[
 |\Gamma(-\chi_k)|\asymp |\chi_k|^{-1/2}\e^{-\pi|\chi_k|/2}.
\]
Since $|\chi_1|=2\pi/L\approx9.0647$, the gamma factor alone contributes an exponential
suppression near $\e^{-14.24}$.  The zeta factor is only of polynomial size.  This explains
the extraordinary numerical smallness of the periodic modes; their nonvanishing follows
separately from the zero-free line $\Re s=1$ for the Riemann zeta function.

\subsection{The unperturbed and perturbed Lambert equations}
The familiar lower-Lambert phase solves
\[
 x2^{u_0}=u_0,
 \qquad
 u_0=-L^{-1}W_{-1}(-Lx).
\]
The exact phase replaces $u_0$ on the right by $u+C(u)$.  The map
\eqref{eq:W-map} is not an ad hoc iteration: it solves the Lambert part exactly while
freezing only the bounded periodic perturbation.  Its derivative
\eqref{eq:W-derivative} is therefore of order $1/u$, explaining the rapid numerical
convergence seen in the table.

\section{Recursive source inventory}
\label{app:inventory}

\subsection{Living documents: five files}
\begin{enumerate}
\item \path{FabiusFunction_Mathematical_Glossary/FabiusFunction_Mathematical_Glossary.tex}
\item \path{Fabius_Function_and_Rvachev_Up/Fabius_Function_and_Rvachev_Up.tex}
\item \path{Non_Elementarity_of_the_Fabius_Function/Non_Elementarity_of_the_Fabius_Function.tex}
\item \path{fabius_lean_walkthrough/fabius_lean_walkthrough.tex}
\item \path{non-formalized-research-frontiers/non-formalized-research-frontiers.tex}
\end{enumerate}

\subsection{Active frontier drafts: thirteen files}
\begin{itemize}
\item Atlas file: \path{non-formalized-research-frontiers/drafts/Thue_Morse_Formula_Atlas/Thue_Morse_Formula_Atlas.tex}.
\item Twelve one-file directories under
\path{non-formalized-research-frontiers/drafts/rvachev_up_fourier_decay/}:
\begin{itemize}
\item \path{Rvachev_Up_Fourier_Decay-2/}
\item \path{Rvachev_Up_Fourier_Decay-3/}
\item \path{Rvachev_Up_Fourier_Decay_Comparative_Audit/}
\item \path{Rvachev_Up_Fourier_Decay_Gentle_Guide/}
\item \path{Rvachev_Up_Fourier_Decay_Second_Wave_Audit/}
\item \path{asymptotic-decay-rate-of-an-infinite-product-of-sinc-functions/}
\item \path{rvachev_up_fourier_decay-1/}
\item \path{rvachev_up_fourier_decay-4/}
\item \path{rvachev_up_fourier_decay-5/}
\item \path{rvachev_up_fourier_decay-6/}
\item \path{rvachev_up_fourier_decay-7/}
\item \path{rvachev_up_fourier_decay-8/}
\end{itemize}
\end{itemize}
Together these are one Atlas file plus twelve Fourier-decay files.

\subsection{Frozen archive: thirty-two files}
The archive contains one historical \texttt{.tex} per entry.  Its README states that
byte-identical duplicates were pruned and that the retained entries are frozen first or
special historical versions.  The entries are:
\begin{longtable}{@{}p{0.34\textwidth}p{0.60\textwidth}@{}}
\toprule
Archive group & Entry\\
\midrule
\endfirsthead
\toprule
Archive group & Entry\\
\midrule
\endhead
Early notes & \path{Fabius Asymptotic}; \path{K-fold summation over the signed Thue-Morse sequence}\\
Standalone studies & \path{Non_Elementarity_of_the_Fabius_Function};
\path{Small_Argument_Asymptotics}; \path{fabius-inverse-dyadic-closed-form}\\
Primary exposition & \path{Fabius_Function_and_Rvachev_Up};
\path{Fabius_Function_and_Rvachev_Up-2};
\path{Fabius_Function_and_Rvachev_Up-3};
\path{Fabius_and_Rvachev_Up-2};
\path{Fabius_and_Rvachev_up};
\path{Fabius_function_and_up_function}\\
Primary supplements & \path{Fabius_Function_Missing_Parts};
\path{Fabius_Function_Missing_Parts-2}\\
$q$-calculus & \path{Fabius_q_Special_Functions}; \path{fabius-q-special-functions}\\
Lean walkthrough & \path{Fabius_Function_Lean_Walkthrough}; \path{fabius_lean_walkthrough}\\
Glossary & \path{FabiusFunction_Mathematical_Glossary};
\path{FabiusFunction_Mathematical_Glossary-2}\\
Research frontiers & \path{Fabius_Dyadic_Asymptotic_Bridge};
\path{Fabius_Dyadic_Formulae_and_Alternative_Representations};
\path{Fabius_Dyadic_Formulae_to_Asymptotics};
\path{Fabius_Dyadic_q_Connections};
\path{Fabius_Thue_Morse_Convergence_Rate};
\path{Fabius_Thue_Morse_Convergence_Rates};
\path{Primary_Exposition_Gap_Register};
\path{Repeated_Integration_and_Rvachev_Up};
\path{Rvachev_Up_from_Repeated_Integration};
\path{non-formalized-research-frontiers};
\path{Thue_Morse_Formula_Atlas-1};
\path{Thue_Morse_Formula_Atlas-2};
\path{Thue_Morse_Formula_Atlas-3}\\
\bottomrule
\end{longtable}

\subsection{Source-paper transcriptions: seven files}
\begin{enumerate}
\item \path{AIF_2017__67_2_637_0/AIF_2017__67_2_637_0-corrected.tex}
\item \path{AIF_2017__67_2_637_0/AIF_2017__67_2_637_0.tex}
\item \path{Lacunary_products/Lacunary_products.tex}
\item \path{09-Function/09-Function.tex}
\item \path{157-Arithmetic-v3/157-Arithmetic-v3.tex}
\item \path{document/document.tex}
\item \path{ubiq15/ubiq15-corrected.tex}
\end{enumerate}

\section{Notation crosswalk}
\begin{center}
\begin{tabularx}{\textwidth}{@{}lX@{}}
\toprule
Symbol & Meaning\\
\midrule
$F$ & bounded Fabius distribution function\\
$\Up$ & Rvachev's even compactly supported density\\
$X$ & $\sum_{j\ge1}2^{-j}U_j$, supported on $[0,1]$\\
$Y$ & $2X-1$, with density $\Up$\\
$\Phi$ & $\widehat\Up$, the binary infinite sinc product\\
$\varepsilon_n$ & signed Thue--Morse value $(-1)^{\wt(n)}$\\
$\Theta(q)$ & Thue--Morse Mahler product $\sum\varepsilon_nq^n$\\
$B,\Lambda$ & negative Laplace transform of $X$ and its logarithm\\
$T,\mathcal E$ & forward logarithmic tail and positive endpoint correction\\
$P$ & smooth one-periodic term in the exact Laplace decomposition\\
$R$ & $sT'(s)$, the exponentially small derivative correction\\
$A$ & $-s\Lambda'(s)=s\E_sX$\\
$\Pi$ & $-1/2-P'/L$, the periodic part of $A-\log_2s$\\
$C$ & $\Pi(u)+R(2^u)$, the exact bounded saddle perturbation\\
$\mu(n)$ & integer-zero multiplicity $\ord_{z=n}\Phi=\vtwo(n)+1$\\
\bottomrule
\end{tabularx}
\end{center}

\begin{thebibliography}{99}

\bibitem{proveit-primary}
V.~Reshetnikov,
\emph{The Fabius Function and Rvachev's Up-Function: Exact arithmetic, probability,
Fourier analysis, and corrected sharp asymptotics},
\texttt{ProveIt} repository, living exposition, snapshot 27 August 2026.
\url{https://github.com/VladimirReshetnikov/ProveIt/tree/main/Analysis/FabiusFunction/docs/Fabius_Function_and_Rvachev_Up}

\bibitem{proveit-frontier}
V.~Reshetnikov,
\emph{Non-formalized Research Frontiers for the Fabius Function},
\texttt{ProveIt} repository, living frontier synthesis, 2026.
\url{https://github.com/VladimirReshetnikov/ProveIt/tree/main/Analysis/FabiusFunction/docs/non-formalized-research-frontiers}

\bibitem{proveit-atlas}
V.~Reshetnikov,
\emph{Thue--Morse Formula Atlas}, active draft, \texttt{ProveIt} repository, 2026.
\url{https://github.com/VladimirReshetnikov/ProveIt/tree/main/Analysis/FabiusFunction/docs/non-formalized-research-frontiers/drafts/Thue_Morse_Formula_Atlas}

\bibitem{proveit-fourier-audit}
V.~Reshetnikov,
\emph{The Decay Envelope of the Rvachev Sinc Product: Comparative Audit},
active draft, \texttt{ProveIt} repository, 2026.
\url{https://github.com/VladimirReshetnikov/ProveIt/tree/main/Analysis/FabiusFunction/docs/non-formalized-research-frontiers/drafts/rvachev_up_fourier_decay/Rvachev_Up_Fourier_Decay_Comparative_Audit}

\bibitem{proveit-fourier-wave2}
V.~Reshetnikov,
\emph{The Decay Envelope of the Rvachev Sinc Product: Second-Wave Audit},
active draft, \texttt{ProveIt} repository, 27 August 2026.
\url{https://github.com/VladimirReshetnikov/ProveIt/tree/main/Analysis/FabiusFunction/docs/non-formalized-research-frontiers/drafts/rvachev_up_fourier_decay/Rvachev_Up_Fourier_Decay_Second_Wave_Audit}

\bibitem{arias-function}
J.~Arias de Reyna,
\emph{An infinitely differentiable function with compact support: Definition and
properties}, English translation of the 1982 article; transcription included in the
\texttt{ProveIt} documentation corpus.

\bibitem{arias-arithmetic}
J.~Arias de Reyna,
\emph{Arithmetic of the Fabius function}, transcription included in the
\texttt{ProveIt} documentation corpus.

\bibitem{allouche-shallit}
J.-P.~Allouche and J.~Shallit,
\emph{The ubiquitous Prouhet--Thue--Morse sequence}, source-paper transcription included
in the \texttt{ProveIt} corpus.

\bibitem{aistleitner}
C.~Aistleitner, M.~Hofer, and G.~Larcher,
\emph{On evil Kronecker sequences and lacunary trigonometric products},
Annales de l'Institut Fourier \textbf{67} (2017), 637--687; corrected transcription included
in the \texttt{ProveIt} corpus.

\bibitem{corless}
R.~M.~Corless, G.~H.~Gonnet, D.~E.~G.~Hare, D.~J.~Jeffrey, and D.~E.~Knuth,
\emph{On the Lambert $W$ function},
Advances in Computational Mathematics \textbf{5} (1996), 329--359.

\end{thebibliography}

\end{document}
